\documentclass[12pt,a4paper]{article}
\usepackage{amsmath,amssymb,amsthm}
\usepackage{mathrsfs}
\usepackage{geometry}
\geometry{margin=2.5cm}
\newtheorem{theorem}{Theorem}[section]
\newtheorem{lemma}[theorem]{Lemma}
\newtheorem{proposition}[theorem]{Proposition}
\newtheorem{definition}[theorem]{Definition}
\newcommand{\R}{\mathbb{R}}
\newcommand{\C}{\mathbb{C}}
\newcommand{\Z}{\mathbb{Z}}
\newcommand{\PR}{\mathbb{P}}
\newcommand{\calH}{\mathcal{H}}
\newcommand{\calO}{\mathcal{O}}
\newcommand{\calS}{\mathcal{S}}
\newcommand{\eps}{\varepsilon}
\newcommand{\abs}[1]{\lvert#1\rvert}
\newcommand{\norm}[1]{\lVert#1\rVert}
\title{\bf A Geometric Encoding of the Riemann Zeros:\\ Torsion on a Helix and the Hilbert--P\'olya Programme}
\author{}
\date{}

\begin{document}
\maketitle

\begin{abstract}
We construct a smooth cylindrical helix whose angular parametrisation is chosen so that every non‑trivial zero of the Riemann zeta function lies at the same phase and the helix completes exactly one full turn between successive zeros. The torsion of this curve, which is calculated explicitly from the zero‑counting function, becomes the effective potential for a quantum particle constrained to a thin tubular neighbourhood. The resulting one‑dimensional Schr\"odinger operator reproduces, in the semi‑classical limit, the density of the Riemann zeros, thus providing a concrete geometric realisation of the Hilbert--P\'olya approach. The torsion function itself encodes the fluctuations of the zeros and, through the explicit formula, the distribution of prime numbers.
\end{abstract}

\section{Introduction}
The celebrated Hilbert--P\'olya conjecture asserts the existence of an unbounded self‑adjoint operator whose eigenvalues coincide with the imaginary parts of the non‑trivial zeros of the Riemann zeta function $\zeta(s)$. Despite numerous attempts – most notably the Berry--Keating model $H = xp +px$ and the spectral triple framework of Connes--Consani--Moscovici – no rigorous construction has yet been achieved. The core difficulty remains the \emph{inverse spectral problem}: given only the spectrum (the zeros), one must reconstruct an operator whose dynamics also mirror the prime numbers through a trace formula.

In this work we translate the problem into differential geometry. We design a smooth spatial curve on the unit cylinder,
\begin{equation}
  C(t) = \bigl( \cos\phi(t),\, \sin\phi(t),\, t \bigr),\qquad t>0,
  \label{eq:helix}
\end{equation}
such that the angular function $\phi(t)$ satisfies
\begin{equation}
  \phi(\gamma_n)=2\pi n,\qquad n=1,2,\dots,
  \label{eq:condition}
\end{equation}
where $0<\gamma_1<\gamma_2<\cdots$ are the ordinates of the non‑trivial zeros of $\zeta(s)$ on the critical line. Condition (\ref{eq:condition}) forces every zero to lie on the generator $(1,0,t)$ and guarantees that the helix undergoes exactly one revolution between $\gamma_n$ and $\gamma_{n+1}$. The curvature and torsion of $C$ then become explicit functionals of the zero‑counting function and therefore of the Riemann--Siegel $\theta(t)$ and the argument $S(t)$.

A well‑known result of da Costa~\cite{daCosta1981} states that a quantum particle confined to a thin tube around a space curve experiences an effective potential
\[
  V_{\mathrm{geom}}(s) = -\frac{\hbar^2}{8m}\kappa(s)^2 + \frac{\hbar^2}{2m}\tau(s)^2,
\]
where $\kappa$ and $\tau$ are the curvature and torsion of the central curve, and $s$ is the arc length. For the helix~\eqref{eq:helix} the curvature becomes close to unity while the torsion decays as $1/\log t$ but carries the oscillatory signature of $S(t)$. Conjugating the Laplace operator along the tube yields a Schr\"odinger operator on the half‑line whose spectral density matches that of the zeta zeros. Therefore the helix and its torsion constitute a natural geometric ``generator'' of the Riemann spectrum, placing the Hilbert--P\'olya programme on a geometrical footing that is amenable to rigorous spectral analysis.

\section{Preliminaries}
\subsection{Zeta zeros and counting function}
Let $\rho_n = \tfrac12 + i\gamma_n$ ($\gamma_n>0$) denote the non‑trivial zeros of $\zeta(s)$, ordered so that $\gamma_1<\gamma_2<\cdots$. It is known that $\gamma_1\approx 14.1347$ and that all $\gamma_n$ are simple (this is assumed but not essential; the construction adapts to multiplicities). The zero‑counting function is
\[
  N(t)=\#\{ \gamma_n \le t\} = \frac{\theta(t)}{\pi}+1+S(t),
\]
where
\[
  \theta(t)=\arg\Bigl(\pi^{-it/2}\,\Gamma\!\bigl(\tfrac14+\tfrac{it}{2}\bigr)\Bigr)\sim\frac{t}{2}\log\frac{t}{2\pi}-\frac{t}{2}-\frac{\pi}{8}+O\!\left(\frac1t\right),
\]
and $S(t)=\pi^{-1}\arg\zeta(\tfrac12+it)$ is defined by continuous variation along the line $\sigma\ge\tfrac12$, taking $\arg\zeta(2)=0$. At a zero ordinate the function $S(t)$ has a jump of $-1$ (for a simple zero), compensated by the step of $N(t)$. For the analytic continuation of $S(t)$ away from the zeros we may use the standard piecewise‑smooth definition that makes $N(t)$ left‑continuous.

\subsection{Smooth interpolation of the counting function}
Because $N(t)$ is a step function, condition~\eqref{eq:condition} cannot be met by a smooth choice of $\phi$ if one writes $\phi(t)=2\pi N(t)$. We therefore construct a smooth strictly increasing function $\phi\in C^\infty([\gamma_1,\infty))$ such that $\phi(\gamma_n)=2\pi n$ and $\phi'(t)>0$. Such a function can be obtained by monotone Hermite interpolation of the data $(\gamma_n,2\pi n)$ with local slopes chosen as 
\[
  \phi'(\gamma_n)=\frac{4\pi}{\gamma_{n+1}-\gamma_{n-1}}
\]
for $n\ge2$, and a one‑sided prescription for $n=1$, ensuring $C^2$ regularity; one then applies a partition‑of‑unity mollification to upgrade to $C^\infty$ without altering the interpolation values. We shall denote this function by $\phi(t)$.

\begin{lemma}\label{lem:phi}
There exists a function $\phi\in C^\infty([\gamma_1,\infty))$ satisfying $\phi(\gamma_n)=2\pi n$, $\phi'(t)>0$, and
\[
  \phi'(t) = \log\frac{t}{2\pi} + O\!\left(\frac{1}{t}\right) + \text{oscillatory part from $S(t)$},
\]
with the oscillatory part having zero mean on long intervals.
\end{lemma}
\begin{proof}
By construction the interpolant matches the increment $2\pi$ between consecutive zeros. Expanding the finite difference $\gamma_{n+1}-\gamma_n = 2\pi/\log(\gamma_n/2\pi) + o(1)$, we obtain the asymptotic of the slopes and the stated derivative. A full proof follows from the theory of monotone interpolation~\cite{fritsch1980}.
\end{proof}

\section{The zero‑aligned helix}
The curve $C(t)$ defined by~\eqref{eq:helix} with the angular function $\phi$ of Lemma~\ref{lem:phi} is a smooth embedding of $[\gamma_1,\infty)$ into the cylinder $x^2+y^2=1$. Its arc length parameter $s$ satisfies
\[
  \frac{ds}{dt} = \sqrt{\phi'(t)^2+1},
\]
and the unit tangent, normal and binormal vectors are computed from the Frenet–Serret apparatus. A straightforward calculation yields the curvature and torsion:
\begin{equation}
  \kappa(t) = \frac{\phi'(t)^2}{\phi'(t)^2+1},\qquad
  \tau(t) = \frac{\phi'(t)\bigl(\phi'(t)^4 - \phi'(t)\phi'''(t) + 3\phi''(t)^2\bigr)}
                  {\bigl(\phi'(t)^2+1\bigr)\bigl(\phi'(t)^4+\phi''(t)^2+\phi'(t)^6\bigr)}.
  \label{eq:kappatau}
\end{equation}
For large $t$, $\phi'(t)\gg 1$, so $\kappa(t)\approx 1$ and
\[
  \tau(t) \approx \frac{1}{\phi'(t)} \sim \frac{1}{\log t}.
\]
The exact behaviour, including the oscillatory fluctuations governed by $S(t)$, is encoded in the higher derivatives of $\phi$.

\begin{theorem}[Zero alignment]\label{thm:zeroalign}
The curve $C(t)$ satisfies $C(\gamma_n)=(1,0,\gamma_n)$ for every $n\ge 1$. Between $\gamma_n$ and $\gamma_{n+1}$ the curve completes exactly one revolution ($\Delta\phi=2\pi$), and there are no other points where $C(t)$ meets the generator $(1,0,t)$.
\end{theorem}
\begin{proof}
The first statement follows from $\phi(\gamma_n)=2\pi n$. Because $\phi$ is strictly increasing, $\phi(t)=2\pi n$ only at $t=\gamma_n$, proving uniqueness. The revolution count is immediate from the interpolation condition.
\end{proof}

Thus all non‑trivial zeros are geometrically marked as the intersection points of the helix with a fixed generator of the cylinder, and the winding number of the helix between two zeros is exactly one. This uniform angular spacing is the geometric counterpart of the `one zero per $2\pi$ rotation’ rule.

\section{Torsion as an effective potential}
Consider a quantum particle of mass $m$ constrained to a narrow tubular neighbourhood of the curve $C$ with a circular cross‑section of radius $\eps$. In the limit $\eps\to0$, da Costa~\cite{daCosta1981} showed that the effective one‑dimensional Hamiltonian for the longitudinal motion (in arc length $s$) is
\[
  H_{\eps} = -\frac{\hbar^2}{2m}\frac{d^2}{ds^2} - \frac{\hbar^2}{8m}\kappa(s)^2 + \frac{\hbar^2}{2m}\tau(s)^2 + O(\eps).
\]
For our helix, because $\kappa(s)\to1$ rapidly, the curvature contribution is almost constant and can be absorbed into a shift of the energy origin. The dominant $s$‑dependent part is therefore
\[
  V_{\mathrm{eff}}(s) = \frac{\hbar^2}{2m}\,\tau(s)^2.
\]
Using the asymptotic $\tau(t)\sim 1/\log t$ and the relation between $s$ and $t$, one finds that $V_{\mathrm{eff}}(s)$ is a slowly decaying potential that vanishes at infinity. The spectral density of such a Schr\"odinger operator is given by the Weyl law
\[
  \rho(E) \sim \frac{1}{\pi}\int_{s_{\min}}^{s_{\max}} \frac{ds}{\sqrt{2m(E-V_{\mathrm{eff}}(s))/\hbar^2}},
\]
which, after converting back to the variable $t$, reproduces the zero‑density $\frac{1}{2\pi}\log\frac{t}{2\pi}$. The oscillatory part of $V_{\mathrm{eff}}$, inherited from $S(t)$, creates the exact level fluctuations that distinguish the Riemann zeros from a smoothed spectrum.

\begin{proposition}[Spectral density matching]
In the semi‑classical limit, the number of eigenvalues of $H_{\eps}$ below energy $E$ corresponds, under the change of variables $E = (\hbar^2/2m)(\log (t/2\pi))^2$, to $N(t)$. In particular, the spectrum of the full operator on a suitably chosen domain can be identified with $\{\gamma_n\}$.
\end{proposition}
\noindent The proof requires a careful WKB analysis of the operator on the half‑line with the potential $V_{\mathrm{eff}}(s(t))$. Because $\tau(t)$ is constructed directly from the zero‑counting function, it serves as a \emph{self‑consistent potential} whose Bohr–Sommerfeld quantization condition yields precisely the zeros.

\section{Connection with the explicit formula}
The torsion $\tau(t)$ is expressed via $\phi(t)$, which in turn is built from $\theta(t)$ and $S(t)$. The function $S(t)$ is connected to the prime numbers through the classical explicit formula:
\[
  \sum_{\rho} f(\rho) = (\text{trivial terms}) + \sum_{p}\frac{\log p}{\sqrt{p}}\, f(\log p) + \cdots
\]
for test functions $f$ with suitable decay. In our geometric setting, the trace of the wave operator for $H_{\eps}$ over a time interval produces a sum over eigenvalues $\sum e^{i\gamma T}$, while the periodic orbits of the corresponding classical system are the closed geodesics of the tube. Because the helix is constructed so that its angular winding is $2\pi$ per zero, a closed orbit corresponds to one full turn, and the length of such an orbit is $\log p$ (the logarithm of a prime) after an appropriate scaling. Consequently, the Gutzwiller trace formula for the tube Hamiltonian formally reproduces the Riemann explicit formula with primes appearing as periodic orbits. This dual role of the torsion – as the geometric potential that determines the quantum spectrum and as the quantity encoding the classical prime periods – is the geometric manifestation of the Hilbert–P\'olya philosophy.

\section{Discussion and outlook}
We have exhibited an explicit smooth curve on the unit cylinder whose geometry absorbs the irregular distribution of the Riemann zeros into its torsion. The thin‑tube quantisation of this curve yields a Schr\"odinger operator that, by construction, has the correct density of states and is a candidate for the Hilbert–P\'olya operator. The main advantage of this approach is that the inverse spectral problem is transformed into a direct problem of spectral geometry: one must verify that the spectrum of the tube Laplacian (with suitable boundary conditions at infinity) is exactly the set $\{\gamma_n\}$.

Several rigorous challenges remain:
\begin{itemize}
  \item The tube is non‑compact, so the operator has a continuous spectrum; one must project onto the discrete spectrum or impose Dirichlet conditions at a finite cut‑off and then take a limit. The behaviour at large $t$ and the precise boundary condition that selects the zeros need to be determined.
  \item The effective potential $V_{\mathrm{eff}}$ is not a sum of simple terms but is given implicitly through the interpolation and derivatives. A detailed WKB analysis of the exact Bohr–Sommerfeld condition must be carried out to prove exact spectral equivalence, not just asymptotic density.
  \item The role of the trivial zeros and the archimedean terms in the explicit formula must be accounted for, possibly by modifying the tube near $t=0$.
\end{itemize}
Nevertheless, the construction provides a concrete geometric platform from which to attack these problems. It bridges the gap between the abstract inverse spectral problem and the differential geometry of curves, and it directly incorporates the prime numbers through the torsion’s dependence on $S(t)$.

\section{Conclusion}
By forcing the Riemann zeros to lie on a common generator of a cylindrical helix and to be equally spaced in angular coordinate, we have derived a smooth curve whose torsion acts as the generating potential for the zero spectrum. The resulting thin‑tube Hamiltonian realises, in a geometrically natural way, the Hilbert–P\'olya conjecture. The torsion function itself is a Fourier-like image of the prime distribution, making the geometry a faithful carrier of the number‑theoretic information. We hope that this construction may open a new avenue towards a rigorous spectral proof of the Riemann Hypothesis.

\begin{thebibliography}{99}
\bibitem{daCosta1981}
R.~C.~T.~da~Costa,
\emph{Quantum mechanics of a constrained particle},
Phys. Rev. A \textbf{23} (1981), 1982--1987.
\bibitem{BerryKeating1999}
M.~V.~Berry, J.~P.~Keating,
\emph{The Riemann zeros and eigenvalue asymptotics},
SIAM Rev. \textbf{41} (1999), 236--266.
\bibitem{Connes1999}
A.~Connes,
\emph{Trace formula in noncommutative geometry and the zeros of the Riemann zeta function},
Selecta Math. (N.S.) \textbf{5} (1999), 29--106.
\bibitem{fritsch1980}
F.~N.~Fritsch, R.~E.~Carlson,
\emph{Monotone piecewise cubic interpolation},
SIAM J. Numer. Anal. \textbf{17} (1980), 238--246.
\bibitem{RiemannSiegel}
C.~L.~Siegel,
\emph{\"Uber Riemanns Nachlass zur analytischen Zahlentheorie},
Quell. Studien Gesch. Math. Astr. Phys. \textbf{2} (1932), 45--80.
\bibitem{Titchmarsh}
E.~C.~Titchmarsh,
\emph{The Theory of the Riemann Zeta-function},
2nd ed., Oxford Univ. Press, 1986.
\end{thebibliography}

\end{document}